\section{Introduction}\label{sec:intro}

Antisemitic incidents do not arrive at a steady rate. They come in surges,
and surges often follow salient external events. In the United Kingdom, the
Community Security Trust (CST) recorded an average of five antisemitic
incidents a day between 1~January and 6~October 2023, and 31 a day from
7~October, the day of the Hamas attack on Israel, to the end of the year. Its
daily count peaked at 80 on 11~October, the highest number ever reported to
CST on a single day \citep{cst2024incidents}. On 2~October 2025, the day of
a terrorist attack on the Heaton Park synagogue in Manchester, and on the
following day, CST logged 40 incidents each, its highest daily totals of the
year. Of those 80 incidents, 42 were antisemitic reactions to the attack
itself, most of them online \citep{cst2026incidents}. In New York City, an analysis of police
records found the rate of anti-Jewish hate crimes in the first year of the
Israel--Hamas war to be about twice that of the preceding years of the study
period \citep{rosenbaum2025}. Such patterns are not peculiar to antisemitism.
In the United States, contentious trial verdicts and lethal terrorist attacks
have preceded spikes in racially or religiously motivated hate crimes
\citep{king2013}.

A surge raises three questions for the communities affected and for those
who plan their security. \emph{How much of it is a direct reaction to the
event?} \emph{How much is incidents begetting incidents}, through imitation,
emboldenment or repeated targeting? And \emph{does online hostility spill
over into offline incidents?} The answers call for different responses. A
direct reaction is concentrated in the days after the event and can be
anticipated from a calendar of events. Contagion means that each incident
prevented also prevents some of the incidents it would have triggered, and
that a surge can outlast its trigger. A spillover from online to offline
streams would make online monitoring a source of early warning. This paper
develops the statistical framework needed to ask these questions of incident
data, to answer them where the data allow, and to say plainly where they do
not.

\IfFileExists{figures/fig-branching.tex}{% =====================================================================
%  Figure: the cluster (immigration--birth) representation (LEAD).
%  Schematic -- event times are hand-placed, NOT data or a simulation.
%  Panel (b) is computed exactly in TikZ from the exponential-kernel
%  intensity, using the Markov recursion between event times.
%  Requires the house colours defined in preamble.tex.
% =====================================================================
\begin{figure}[t]
\centering
\begin{tikzpicture}[x=0.952cm,y=1cm,
    ev/.style={circle,inner sep=0pt,minimum size=5.2pt},
    bg/.style={ev,fill=LabNavyLight,draw=white,line width=0.5pt},
    sh/.style={ev,fill=LabAccent,draw=white,line width=0.5pt},
    of/.style={ev,fill=LabNavy,draw=white,line width=0.5pt},
    kid/.style={LabNavyMid!75,line width=0.55pt,
                -{Stealth[length=3.2pt,width=2.6pt]},shorten >=1.6pt,shorten <=1.4pt},
    genlab/.style={font=\scriptsize,text=LabGrey,anchor=east,inner sep=1pt},
    panel/.style={font=\footnotesize\bfseries,text=LabNavy,anchor=south west,inner sep=0pt},
    note/.style={font=\scriptsize\itshape,text=LabGrey,inner sep=1pt},
    leg/.style={font=\scriptsize,text=LabInk,anchor=west,inner sep=1pt},
    axis/.style={LabNavy,line width=0.5pt},
  ]
  % ---------------- parameters (schematic) ----------------
  \def\gy{0.68}       % vertical spacing of generations in panel (a)
  \def\tshock{5.6}    % time of the salient external event
  \def\tend{15.3}     % right end of the time axis

  % =========================== panel (a) ===========================
  \node[panel] at (-1.05,3*\gy+0.55) {(a)\enspace Who triggered whom
    \textnormal{\textcolor{LabGrey}{(latent, unobserved)}}};
  \foreach \g in {0,1,2,3}{
    \draw[LabRule,line width=0.3pt,dash pattern=on 1pt off 1.5pt] (0,\g*\gy) -- (\tend,\g*\gy);
    \node[genlab] at (-0.12,\g*\gy) {gen.\,\g};}

  % a highlighted cluster (one immigrant and all its descendants)
  \begin{scope}[on background layer]
    \fill[LabTint,rounded corners=4pt] (0.38,-0.2) rectangle (2.55,2*\gy+0.22);
  \end{scope}
  \node[note,anchor=south] at (1.47,2*\gy+0.22) {one cluster};

  % background immigrants (generation 0)
  \node[bg] (b1) at (0.6,0) {};
  \node[bg] (b2) at (3.4,0) {};
  \node[bg] (b3) at (10.6,0) {};
  \node[bg] (b4) at (13.6,0) {};
  % shock-triggered immigrants (generation 0)
  \node[sh] (s1) at (5.75,0) {};
  \node[sh] (s2) at (6.05,0) {};
  \node[sh] (s3) at (6.4,0) {};
  \node[sh] (s4) at (6.95,0) {};
  \node[sh] (s5) at (7.7,0) {};
  % offspring
  \node[of] (o1) at (1.3,\gy) {};
  \node[of] (o2) at (2.2,2*\gy) {};
  \node[of] (o3) at (4.3,\gy) {};
  \node[of] (o4) at (14.5,\gy) {};
  \node[of] (c1) at (6.5,\gy) {};
  \node[of] (c2) at (7.3,\gy) {};
  \node[of] (c3) at (6.85,\gy) {};
  \node[of] (c4) at (7.05,2*\gy) {};
  \node[of] (c5) at (8.3,2*\gy) {};
  \node[of] (c6) at (8.9,3*\gy) {};
  \node[of] (c9) at (8.6,2*\gy) {};
  \node[of] (c8) at (9.4,\gy) {};
  \node[of] (c10) at (10.1,2*\gy) {};
  \node[of] (c11) at (11.4,\gy) {};
  \node[of] (c12) at (12.3,2*\gy) {};
  % parent -> child
  \foreach \p/\c in {b1/o1,o1/o2,b2/o3,b4/o4,s1/c1,s1/c2,s2/c3,c1/c4,c1/c5,
                     c4/c6,c2/c9,s4/c8,c8/c10,s5/c11,c11/c12}{
    \draw[kid] (\p) to[out=72,in=208,looseness=0.85] (\c);}

  % legend (one row, right-aligned with the time axis)
  \begin{scope}[node distance=1.5pt]
    \node[leg,anchor=east] (l3) at (\tend+0.25,3*\gy+0.66) {offspring};
    \node[of,left=of l3] (d3) {};
    \node[leg,anchor=east] (l2) at ([xshift=-9pt]d3.west) {shock-triggered immigrant};
    \node[sh,left=of l2] (d2) {};
    \node[leg,anchor=east] (l1) at ([xshift=-9pt]d2.west) {background immigrant};
    \node[bg,left=of l1] {};
  \end{scope}

  % =========================== panel (b) ===========================
  \begin{scope}[yshift=-3.75cm,y=0.45cm]
    \node[panel] at (-1.05,5.6) {(b)\enspace What is recorded: event times, and
      the conditional intensity $\lambda(t)$ they imply under the model};
    % model parameters (schematic): background, kernel, shock response
    \def\kmu{0.6}\def\kalpha{0.55}\def\kbeta{1.0}\def\kgam{0.8}\def\kjump{4.0}
    \xdef\tp{0}\xdef\Hs{0}\xdef\Ss{0}
    % breakpoints: time/kind, kind 1 = incident, 2 = shock, 0 = end of axis
    \foreach \t/\k in {0.6/1,1.3/1,2.2/1,3.4/1,4.3/1,5.6/2,5.75/1,6.05/1,6.4/1,
                       6.5/1,6.85/1,6.95/1,7.05/1,7.3/1,7.7/1,8.3/1,8.6/1,8.9/1,
                       9.4/1,10.1/1,10.6/1,11.4/1,12.3/1,13.6/1,14.5/1,\tend/0}{
      \begin{scope}[on background layer]
        \filldraw[LabNavyLight!45,line width=0.25pt] (\tp,0) rectangle (\t,\kmu);
        \filldraw[LabAccent!55,line width=0.25pt]
          plot[domain=\tp:\t,samples=24] (\x,{\kmu+\Ss*exp(-\kgam*(\x-\tp))})
          -- (\t,\kmu) -- (\tp,\kmu) -- cycle;
        \filldraw[LabNavyMid!80,line width=0.25pt]
          plot[domain=\tp:\t,samples=24]
            (\x,{\kmu+\Ss*exp(-\kgam*(\x-\tp))+\Hs*exp(-\kbeta*(\x-\tp))})
          -- plot[domain=\t:\tp,samples=24] (\x,{\kmu+\Ss*exp(-\kgam*(\x-\tp))})
          -- cycle;
      \end{scope}
      \draw[LabNavy,line width=0.7pt,line join=round]
        plot[domain=\tp:\t,samples=24]
          (\x,{\kmu+\Ss*exp(-\kgam*(\x-\tp))+\Hs*exp(-\kbeta*(\x-\tp))});
      \pgfmathsetmacro{\Hn}{\Hs*exp(-\kbeta*(\t-\tp))+(\k==1)*\kalpha*\kbeta}
      \pgfmathsetmacro{\Sn}{\Ss*exp(-\kgam*(\t-\tp))+(\k==2)*\kjump}
      \pgfmathsetmacro{\Lold}{\kmu+\Ss*exp(-\kgam*(\t-\tp))+\Hs*exp(-\kbeta*(\t-\tp))}
      \pgfmathsetmacro{\Lnew}{\kmu+\Sn+\Hn}
      \ifnum\k>0 \draw[LabNavy,line width=0.7pt] (\t,\Lold) -- (\t,\Lnew);\fi
      \ifnum\k=1 \draw[LabInk,line width=0.6pt] (\t,-0.12) -- (\t,-0.85);\fi
      \xdef\tp{\t}\xdef\Hs{\Hn}\xdef\Ss{\Sn}
    }
    % axes
    \draw[axis,-{Stealth[length=4pt]}] (0,0) -- (\tend+0.25,0)
      node[right,font=\footnotesize,text=LabNavy,inner sep=1pt] {$t$};
    \draw[axis] (0,0) -- (0,5.2);
    \node[genlab,anchor=east] at (-0.1,4.6) {$\lambda(t)$};
    \node[genlab,anchor=east] at (-0.1,\kmu) {$\mu$};
    \node[genlab,anchor=east] at (-0.1,-0.5) {events};
    % layer labels (direct labelling, short leaders)
    \node[font=\scriptsize,text=LabAccent,anchor=east,inner sep=1pt] (lshock)
      at (5.3,3.3) {shock response $\beta_kS_e\,\psi_k(t-\tau_e)$};
    \draw[LabAccent,line width=0.4pt] (lshock.east) -- (6.05,2.0);
    \fill[LabAccent] (6.05,2.0) circle (0.9pt);
    \node[font=\scriptsize,text=LabNavyMid,anchor=west,inner sep=1pt] (lexc)
      at (10.3,3.6) {excitation by earlier incidents $\textstyle\sum_{t_i<t}\varphi(t-t_i)$};
    \draw[LabNavyMid,line width=0.4pt] (lexc.west) -- (9.45,1.45);
    \fill[white] (9.45,1.45) circle (0.9pt);
    \node[font=\scriptsize,text=LabNavy,anchor=west,inner sep=1pt] at (12.35,0.3) {background $\mu$};
  \end{scope}

  % the salient external event, spanning both panels
  \draw[LabAccent,line width=0.7pt,dash pattern=on 2.5pt off 1.5pt]
    (\tshock,-3.75-0.45) -- (\tshock,-3.75+5.0*0.45)
    (\tshock,-0.32) -- (\tshock,3*\gy+0.12);
  \node[font=\scriptsize,text=LabAccent,anchor=south,align=center,inner sep=1pt]
    at (\tshock,3*\gy+0.12) {salient external event $\tau_e$};
\end{tikzpicture}
\caption{\textbf{The cluster representation of a Hawkes process with
exogenous shocks.} \emph{(a)} Every incident is either an \emph{immigrant}
or the \emph{offspring} of an earlier incident. Background immigrants (light
navy) arrive at rate $\mu_k(t)$; shock-triggered immigrants (orange) are the
direct response to a salient external event at $\tau_e$; each incident, of
any kind, independently produces $\Pois(\alpha_{kj})$ offspring at delays
drawn from $g_{kj}$ (navy), which produce offspring in turn
\citep{hawkes1974}. An immigrant and all its descendants form one cluster.
\emph{(b)} The data are the event times alone (ticks). Under the model, the
conditional intensity is the sum of the background, the shock response and
the excitation by past incidents (\cref{eq:th-intensity}). The share of each
layer at an event time is the probability that the event is a background
immigrant, a shock-triggered immigrant or the offspring of an earlier
incident; this is how the EM algorithm reconstructs the unobserved family
trees (\cref{alg:em}). The surge decomposition of \cref{sec:attribution}
classifies incidents instead by the root of their cluster
(\cref{fig:attribution}). Schematic: event times are
hand-placed, and the intensity is computed exactly from them for
illustrative parameters ($\mu=0.6$, $\alpha=0.55$, exponential kernel
with $\beta=1$, and $\beta_kS_e=5$ expected direct incidents with
response $\psi_k(t)=\eta e^{-\eta t}$, $\eta=0.8$).}
\label{fig:branching}
\end{figure}
}{}%
\paragraph{Self-exciting point processes}
The natural model is the self-exciting, or Hawkes, point process
\citep{hawkes1971}, in which every incident raises the short-run rate of
further incidents. Its cluster representation \citep{hawkes1974} makes the
three questions concrete (\cref{fig:branching}). Each incident is either an
\emph{immigrant}, arriving from a background rate or as a direct response to
a salient external event, or the \emph{offspring} of an earlier incident.
The expected numbers of offspring form a branching matrix
$\Amat=[\alpha_{kj}]$ across streams (regions, offence types, online and
offline channels), and the multiplier $\Mmat=(\Imat-\Amat)^{-1}$ counts the
whole cascade that one incident sets off. Self-exciting models are
established in seismology \citep{ogata1988}, crime
\citep{mohler2011,reinhart2018}, conflict \citep{lewis2012}, online popularity
\citep{zhao2015,rizoiu2017} and gun violence \citep{loeffler2018}, and online
content has been found to predict offline hate crime \citep{muller2021,williams2020}.
As far as we know, no study has estimated a multivariate branching matrix for
antisemitic incidents, and this paper does not either. It builds and tests the
tools that such an analysis requires.

\paragraph{Three traps}
Applied naively, the model is easy to fit and easy to over-read. We
formalise three traps.
\begin{enumerate}[(i),leftmargin=2.2em]
  \item \emph{Contagion and clustered shocks look alike, within a stream.} A
  stationary Hawkes process with an exponential kernel has exactly the same mean
  and Bartlett spectrum as a one-parameter family of shot-noise Cox processes, in
  which unrecorded shocks trigger incidents that do not trigger one another
  (\cref{thm:equiv}). For a single stream, the autocorrelation of counts, the
  variance--time curve and the periodogram therefore cannot tell contagion from
  clustered exogenous shocks, at any bin width. This is a dynamic counterpart of
  the old distinction between true and apparent contagion \citep{feller1943}.
  Between streams the picture is different: mutual excitation that dominates
  self-excitation leaves a second-order trace that no Cox process can reproduce
  (\cref{ex:th-mv-cross}).
  \item \emph{A rigid background manufactures contagion.} If salient events
  are missing from the calendar, or the background is too stiff to follow
  slow variation, a Hawkes model with a constant background has a positive
  pseudo-true branching ratio even when there is no contagion at all
  (\cref{thm:spurious}). For second-order estimators the spurious branching
  ratio is explicit, $n^*=1-(1+\kappa)^{-1/2}$, where $\kappa$ measures the
  burstiness of the unrecorded driver. Maximum likelihood need not reproduce
  $n^*$ exactly. In our simulations a single spline background removed the
  artefact of slow rate changes but not that of fast unrecorded shocks, and
  known event times largely removed it where tested (\cref{sec:sim-spurious}).
  For rare large shocks, maximum likelihood settled 0.06--0.07 below $n^*$ even
  with $T=10{,}000$.
  \item \emph{Under-reporting masks contagion.} Most incidents are never
  recorded. That is the subject of the companion paper
  \citep{kornsteinlab2026darkfigure}, which estimates the size of this dark
  figure. If each incident is recorded independently with probability $p$,
  the recorded stream behaves, to second order, like a Hawkes process with
  an attenuated branching ratio $n'$, with $pn<n'<n$ and a closed-form map
  between the two (\cref{thm:thinning}). A cascade passes through unrecorded
  incidents too, so ignoring the dark figure understates contagion. When the
  most productive incidents are also the most likely to be recorded, the
  standard correction overshoots instead (\cref{prop:th-hetthin}).
\end{enumerate}
Two further cautions are quieter. Aggregating incidents into daily counts
also attenuates estimated contagion, by an amount that can be computed
(\cref{prop:est-binning}). And although online-to-offline Granger
non-causality is equivalent to a zero cross-excitation coefficient
(\cref{prop:granger}, after \citealp{eichler2017}), unrecorded common shocks,
classifier error, under-recording and report-time stamps can each distort
that coefficient (\cref{rem:th-granger-caveats}).

\paragraph{What this paper does, and does not, do}
This is a methods paper. We set out the model, prove what can and cannot be
identified, derive the surge decomposition and an alert rule, give
estimation and testing procedures, evaluate some of them on synthetic data
with known ground truth, and pre-register an analysis plan for real data. We
report \emph{no} empirical estimates, and nothing here should be read as a
finding that antisemitic incidents are, or are not, contagious; the figures
quoted above serve only to motivate the questions. \Cref{fig:network} shows
the multivariate model; notation is collected in \cref{tab:notation}.

\IfFileExists{tables/tab-notation.tex}{% =====================================================================
%  Table: notation (LEAD).  Symbols follow SPEC.md and Sections 3--7.
% =====================================================================
\begin{table}[!tp]
\caption{\textbf{Notation.} Symbols used throughout. Bold capitals are
$M\times M$ matrices indexed by (child stream $k$, parent stream $j$).}
\label{tab:notation}
\small
\renewcommand{\arraystretch}{1.08}
\newcommand{\nhead}[1]{\multicolumn{2}{@{}l}{\footnotesize\scshape\bfseries\color{LabNavy}#1}\\[1pt]}
\begin{minipage}[t]{0.485\textwidth}
\begin{tabularx}{\linewidth}[t]{@{}>{$}l<{$}@{\hspace{0.8em}}>{\raggedright\arraybackslash}X@{}}
\toprule
\nhead{Streams and incidents}
k,j\in[M] & stream: offline incidents by region or offence type; online volume \\
t^k_i,\ x^k_i & time and mark of the $i$th stream-$k$ incident on $[0,T]$ \\
N_k(t),\ \Hist_t & counting process; history of all streams and recorded events \\
\lambda_k(t),\ \Lam_k(t) & conditional intensity; compensator \\
\midrule
\nhead{Self-excitation}
\varphi_{kj}=\alpha_{kj}g_{kj} & excitation kernel; $g_{kj}$ a probability density \\
\alpha_{kj} & expected \emph{direct} stream-$k$ offspring of one stream-$j$ incident \\
\Amat=[\alpha_{kj}],\ \spr(\Amat) & branching matrix; spectral radius, $<1$ for stationarity \\
\Mmat=(\Imat-\Amat)^{-1} & contagion multiplier: $\Mmat_{kj}$ is the expected stream-$k$ size of a cascade rooted in stream $j$ \\
\Phik(t),\ \Rker=\sum_{m\ge1}\Phik^{*m} & kernel matrix; resolvent \\
\beta e^{-\beta t} & exponential kernel with decay rate $\beta$ (or $\beta_{kj}$) \\
n,\ \bar\Lambda=\mu/(1-n) & univariate branching ratio; stationary mean rate \\
\bottomrule
\end{tabularx}
\end{minipage}\hfill
\begin{minipage}[t]{0.485\textwidth}
\begin{tabularx}{\linewidth}[t]{@{}>{$}l<{$}@{\hspace{0.8em}}>{\raggedright\arraybackslash}X@{}}
\toprule
\nhead{Background and salient external events}
\mu_k(t)=e^{B(t)^{\top}\theta_k} & background; basis $B$ of $K_b$ functions (spline knot spacing $h_b$) that carries calendar effects \\
e;\ \tau_e,\ S_e & salient external event; its time and size \\
\psi_k,\ \upsilon_k & response kernel ($\int\psi_k=1$, e.g.\ $\eta_ke^{-\eta_kt}$); shock loading \\
\vartheta & full parameter vector, including the background coefficients $\theta_k$ \\
\midrule
\nhead{Recording}
p_k=1/\mathcal D_k & recording probability; $\mathcal D_k$ is the companion paper's dark-figure ratio \\
n'=h_p(n) & apparent branching ratio of the $p$-thinned process \\
d,\ F_k & reporting delay; its distribution \\
e_0,\ e_1 & classifier flag rates for non-antisemitic and antisemitic posts \\
\midrule
\nhead{Attribution and alerting}
B_k,\ D_k,\ C_k & background, direct and contagion parts of $\E\,N_k(\tau_e,\tau_e+H]$ \\
\bm C^{\mathrm{on}},\ \sigma_k(H) & online-channel contagion; its share \\
\gamma=\beta(1-n) & decay rate of a cascade (univariate, exponential kernel) \\
\Delta_{\mathrm a},\ c & alert window; count threshold ($\Delta$ without subscript is a bin width for counts) \\
\bm\chi(t) & excitation state $(\bm x(t),\bm u(t))$ of the alert rule \\
\pi_t,\ \pi^* & surge probability; alert threshold \\
c_{\mathrm{FA}},\ c_{\mathrm{miss}} & costs of a false alarm and of a missed surge \\
\bottomrule
\end{tabularx}
\end{minipage}
\end{table}
}{}%
\paragraph{Contributions}
Our contributions are a synthesis, generalisation and application of
existing ideas to a new and sensitive domain. Classical results are
attributed where they are used; we flag the results that are, to our
knowledge, new in the form given.
\begin{enumerate}[label=\textbf{C\arabic*},ref=C\arabic*,leftmargin=2.4em,itemsep=2pt]
  \item\label{contrib:framework} \textbf{Framework.} A multivariate Hawkes
  process with exogenous shocks for incident streams across regions, offence
  types and online and offline channels (\cref{sec:framework}), and what
  ``contagion'' does and does not mean when it is applied to antisemitic
  incidents (\cref{sec:th-meaning}).
  \item\label{contrib:id} \textbf{Identification.} For a single stream, or
  for streams that do not excite one another, the observational equivalence of
  contagion and clustered shocks is exact, for exponential
  (\cref{thm:equiv}) and completely monotone kernels (\cref{prop:th-cm}). With
  cross-excitation, equivalence requires the excess spectral matrix
  $\bm E_H(\omega)$ to be positive semidefinite at every frequency; this fails
  for excite-only streams and when cross-excitation dominates self-excitation,
  and contagion between streams is then identified at second order
  (\cref{prop:th-mv,ex:th-mv-cross}). Identification is restored by a recorded
  calendar, higher-order information or provenance links
  (\cref{cor:equiv-id}), and a misspecified background yields spurious
  contagion (\cref{thm:spurious}).
  \item\label{contrib:thin} \textbf{Under-reporting.} A closed-form map from
  $n$ to the apparent $n'=h_p(n)$ (\cref{thm:thinning}), whose inversion with
  the companion paper's estimate of $p$ gives a dark-figure correction with a
  variance floor (\cref{cor:th-darkfig}), and the bias of that correction
  under heterogeneous recording (\cref{prop:th-hetthin}).
  \item\label{contrib:attr} \textbf{Surge attribution.} An exact split of the
  expected count after an event into background, direct and contagion
  components (\cref{thm:attribution}), hypoexponential closed forms for
  exponential kernels (\cref{prop:irf}), and a channel decomposition that
  isolates the online-to-offline spillover as a sum over paths
  (\cref{prop:at-channel}); the last two are, to our knowledge, new in this
  form.
  \item\label{contrib:est} \textbf{Estimation and testing.} EM with an exact
  treatment of the observation window (\cref{alg:em,prop:est-mstep}); Bayesian
  shrinkage with an entrywise description of the stationarity region
  (\cref{lem:est-trunc}); an exact formula for the attenuation of contagion by
  binning (\cref{prop:est-binning}); boundary-corrected and sup-type
  likelihood-ratio tests (\cref{sec:est-lrt}); time-rescaling goodness of fit
  (\cref{sec:est-gof}); and placebo calendars (\cref{sec:est-placebo}). The
  placebo rank test is valid whatever the misspecification of the base model,
  but only under the randomisation hypothesis that the real event dates are
  exchangeable with dates drawn from the placebo design
  (\cref{lem:est-placebo}); without it, the placebo distribution is a
  calibration, not a test.
  \item\label{contrib:alert} \textbf{Forecasting and alerting.} When the
  model's state is known, the Bayes-optimal alert rule thresholds the
  predictive probability of a surge at
  $\pi^*=c_{\mathrm{FA}}/(c_{\mathrm{FA}}+c_{\mathrm{miss}})$, and its alert
  region is an upper set in the excitation state (\cref{thm:alert}). Reporting
  delays bias the state downward; a delay-weighted nowcast corrects the bias
  in the mean (\cref{prop:nowcast}). Plugging the nowcast into the rule is not
  the Bayes rule for delayed data. An imputation algorithm that targets it
  (\cref{alg:alert}) is given but was not run.
  \item\label{contrib:sim} \textbf{A simulation study.} Simulations on
  synthetic data with known ground truth evaluate EM, surge attribution, the
  boundary test, the thinning correction and a plug-in alert with a nowcast
  state (\cref{sec:simulation}). Under the model these tools perform as the
  theory says. The study also shows their limits: a single spline background
  did not remove spurious contagion from fast unrecorded shocks, and intervals
  fail when excitation is weak and recording is poor; a 30--50\% error in the
  recording probability sets a floor under the precision of the corrected
  estimate that more data cannot lower. Forecasts that ignore future external
  events under-forecast; forecasts that integrate over the event process are
  calibrated when windows are pooled. What no forecaster can do is time new
  events, which limits recall rather than pooled calibration. On
  Hawkes-generated data the plug-in rule beats two simple baselines, which are
  not the plan's. Not
  evaluated are Bayesian shrinkage, time-rescaling goodness of fit, the
  imputation-based alert of \cref{alg:alert}, the sup-statistic for a decay
  identified only under the alternative (with its Davies approximation and
  bootstrap), nonparametric kernels, the INAR approximation, placebo
  calibration and the CUSUM monitor.
  \item\label{contrib:plan} \textbf{A pre-registered analysis plan} for New
  York City police hate-crime records \citep{nypdhate}, with an event list
  completed blind to incident data, registered block tests, sensitivity
  analyses, placebo shocks and decision rules (\cref{sec:plan}), and an ethics
  protocol (\cref{sec:ethics}).
\end{enumerate}

\paragraph{Related work}
The likelihood theory of self-exciting processes goes back to
\citet{ogata1978} and \citet{ogata1988}, and their stability theory to
\citet{bremaud1996}. EM over the latent branching structure was developed by
\citet{veen2008} and, nonparametrically, by \citet{marsan2008}; book-length
treatments are \citet{daley2003} and \citet{laub2021}. The equivalence of
Hawkes processes and INAR($\infty$) sequences \citep{kirchner2016} underlies
our treatment of binned counts, and the Granger-causal graphs of
\citet{eichler2017} our treatment of online-to-offline dependence. Including
news events in the background lowers estimated endogeneity in financial
markets \citep{rambaldi2015}, the empirical face of \cref{thm:spurious}. Our
nowcasting builds on Bayesian nowcasting in epidemiology
\citep{hohle2014,mcgough2020}.

\paragraph{Roadmap}
\Cref{sec:data} sets out what the data must provide, \cref{sec:framework}
defines the model, and \cref{sec:identification} develops the identification
results. \Cref{sec:attribution} derives the surge decomposition,
\cref{sec:estimation} estimation and testing, and \cref{sec:forecasting}
forecasting and alerting. \Cref{sec:simulation} reports the simulation study,
\cref{sec:plan,sec:ethics} summarise the analysis plan and the ethics
protocol, and \cref{sec:discussion} discusses limitations and extensions. The
appendices contain the proofs (\cref{app:idproofs,app:attrproofs,app:est-proofs}),
supplementary material for \cref{sec:estimation,sec:forecasting}
(\cref{app:est-supp,app:fc-supp}), computational and simulation details
(\cref{app:computation,app:simdetails}), the candidate data sources
(\cref{app:data}), and the full plan and ethics protocol
(\cref{app:plan,app:ethics}).
